\chapter{Periods as functions: Picard--Fuchs equations}\label{ch:pf}

A one-parameter family of K3 surfaces $X_z$ has a holomorphic $2$-form
$\sigma_z$ varying with $z$, and its integrals over a basis of transcendental
cycles satisfy a linear differential equation in $z$, the Picard--Fuchs
equation. Its order is the rank of the transcendental lattice; for the
$M_n$-polarized families that is $3$, and Doran's theorem says the equation is
then always a \emph{symmetric square} of a second-order one. This chapter
explains that circle of ideas and connects it to the Ap\'ery-like sequences
of Zagier and Cooper, which is where the frontier of Chapter~\ref{ch:frontier}
lives.

\section{Variation of Hodge structure and the Picard--Fuchs operator}

Let $\pi\colon\mathcal X\to S$ be a smooth family of K3 surfaces over a curve
$S$, and let $\mathcal H=R^2\pi_*\Z$ be the local system of second cohomology,
with fibres $\Lambda_{K3}$. The Gauss--Manin connection $\nabla$ differentiates
sections. Let $\sigma\in\Gamma(S,\mathcal H\otimes\mathcal O)$ be a family of
holomorphic $2$-forms; then $\sigma,\nabla\sigma,\nabla^2\sigma,\dots$ are
sections of $\mathcal T\otimes\mathcal O$ where $\mathcal T$ is the local
system of transcendental lattices (assuming $\rho$ constant), of rank $r$, and
they become linearly dependent at step $r$:

\begin{definition}
The \emph{Picard--Fuchs operator} of the family is the monic differential
operator $L=\partial_z^r+c_{r-1}(z)\partial_z^{r-1}+\dots+c_0(z)$, with
$c_i\in\C(z)$, such that $L\sigma=0$; equivalently, the period functions
$\int_{\gamma}\sigma_z$ for $\gamma\in\mathcal T^{\vee}$ span its solution
space. Its order is $r=\rk T$.
\end{definition}

In the $\theta$-form, $\theta=z\,d/dz$, the operator is
$L=\sum_i P_i(z)\theta^i$ with polynomial coefficients; this is the form in
which Ap\'ery-like recurrences arise, since $\theta z^n=nz^n$ turns
$L\,\sum a_nz^n=0$ into a recurrence for $(a_n)$.

\begin{theorem}[\Lstatus, Griffiths transversality; regular singularities]
$\nabla\sigma\in\mathcal F^1\otimes\mathcal O$, so $\nabla\sigma\perp\sigma$;
the Picard--Fuchs equation is Fuchsian (regular singular points), with
monodromy the image of $\pi_1(S)$ in $\Orth(T)$ \cite[Ch.~10]{Voisin2002}.
\end{theorem}

\section{Doran's theorem}

\begin{theorem}[\Lstatus, Doran~\cite{Doran2000}, Thm.~5.13]\label{thm:doran}
The Picard--Fuchs equation of a family of $M_n$-polarized K3 surfaces is the
symmetric square of a second-order homogeneous linear Fuchsian ordinary
differential equation.
\end{theorem}

Two proofs are indicated in~\cite{Doran2000}. The first is a general
principle: the period point of an $M_n$-polarized surface lies on the quadric
$\sigma^2=0$ in $\PP(T_\C)\cong\PP^2$, so the three period functions satisfy
a nondegenerate quadratic relation; and a third-order equation whose solutions
satisfy a nondegenerate quadratic relation is the symmetric square of a
second-order one (Doran's Cor.~5.8, which we quote). Explicitly, in the
parametrization $\omega(\tau)=e-n\tau^2f+\tau w$ of Chapter~\ref{ch:torelli},
the three coordinates $(1,-n\tau^2,\tau)$ are the monomials $1,\tau^2,\tau$ in
the two solutions $1,\tau$ of a second-order equation --- which is exactly
what a symmetric square means. The second proof uses the Shioda--Inose
presentation of the surfaces as coming from products of $n$-isogenous
elliptic curves.

\begin{definition}[symmetric square]
For $\Ltwo=\theta^2+p_1\theta+p_0$ with solutions $y_1,y_2$, the
\emph{symmetric square} $\Sym^2\Ltwo$ is the third-order operator annihilating
$y_1^2,y_1y_2,y_2^2$. In $\theta$-form with polynomial coefficients
the coefficients of $\Sym^2\Ltwo$ are explicit polynomials in those of
$\Ltwo$ and their $\theta$-derivatives; the recipe is classical (see e.g.\
\cite[\S4]{AlmkvistVanStraten2021}).
\end{definition}

\begin{warning}\label{warn:converse}
Doran's theorem runs from geometry to analysis: \emph{if} the family is
$M_n$-polarized \emph{then} its operator is a symmetric square. The converse
--- that an order-3 operator which happens to be a symmetric square comes from
an $M_n$-polarized family --- is not a theorem; Doran himself notes in
\cite[\S6]{Doran2000} that the classification of rank-19 lattice polarizations
needed for such a converse is lacking. Chapter~\ref{ch:frontier} lives
precisely in the gap.
\end{warning}

\section{Ap\'ery-like sequences and their operators}

\begin{example}[Ap\'ery's $\zeta(3)$ sequence]
Ap\'ery's proof of the irrationality of $\zeta(3)$ uses
$a_n=\sum_k\binom nk^2\binom{n+k}k^2$, satisfying
\[
  (n+1)^3a_{n+1}-(34n^3+51n^2+27n+5)a_n+n^3a_{n-1}=0 .
\]
Beukers and Peters~\cite{BeukersPeters1984} showed that $\sum a_nz^n$ is a
period of a one-parameter family of K3 surfaces with generic Picard number
$19$ --- the first appearance of Corollary~\ref{cor:rho19} in the wild ---
and Peters--Stienstra identified the family as $M_6$-polarized in today's
language.
\end{example}

\begin{theorem}[\Lstatus, Zagier's sporadic sequences~\cite{Zagier2009}]
Among the second-order recurrences $(n+1)^2u_{n+1}-(an^2+an+b)u_n+cn^2u_{n-1}=0$
with $u_0=1$, $u_{-1}=0$, exactly six ``sporadic'' parameter triples
$(a,b,c)$ (beyond degenerate and Legendrian ones) give integer sequences; each
$\sum u_nz^n$ is a modular form of weight $1$ in a Hauptmodul of a genus-zero
group of level $\le 8$ or so.
\end{theorem}

\begin{theorem}[\Lstatus, Cooper~\cite{Cooper2012}; Gorodetsky~\cite{Gorodetsky2021}]\label{thm:cooper}
There are three further sporadic \emph{third-order} sequences, $s_7,s_{10},s_{18}$,
satisfying recurrences of the form
\[
  (n+1)^3s(n+1) = (2n+1)(an^2+an+b)\,s(n)-n\,(cn^2+d)\,s(n-1),
\]
with $(a,b,c,d)$ equal to $(13,4,-27,3)$, $(6,2,-64,4)$, $(14,6,192,-12)$
respectively, all integer-valued. Their generating functions are weight-2
modular forms in Hauptmoduln of $\Gamma_0(7)^+$, $\Gamma_0(10)^+$, $\Gamma_0(18)^+$.
\end{theorem}

The four-parameter recurrence is Gorodetsky's \emph{template}; the order-3
operator it defines, $\Lthree$, is the object whose symmetric-square structure
is the first result of Chapter~\ref{ch:frontier}. Let us record here the
elementary but central observation that connects the template to
Theorem~\ref{thm:doran}.

\begin{proposition}[\Kstatus; the template is a symmetric square]\label{prop:template}
For every $(a,b,c,d)\in\Q^4$, the $\theta$-form operator $\Lthree$ of the
template above factors as $\Lthree=P_2\cdot\Sym^2(\Ltwo)$ for an explicitly
determined second-order operator $\Ltwo$ whose coefficients are polynomials in
$z$ with coefficients in $\Q[a,b,c,d]$, and the leading coefficient $P_2$ of
$\Ltwo$ satisfies $\theta(P_2)=2P_1$ where $P_1$ is the next coefficient. The
statement is four polynomial identities, each closed by \artifact{ring}.
\end{proposition}

The proof is in \artifact{Agora/Sequences/PartnerOperators.lean}; what it
proves, and what its statement does \emph{not} include, is discussed in
Chapters~\ref{ch:frontier} and~\ref{ch:formal}. We only note here that the
existence of a symmetric-square structure for these three operators is
\emph{consistent with} Doran's theorem, and that consistency is not
confirmation (Warning~\ref{warn:converse}): the template is a symmetric square
for all $(a,b,c,d)$, including the uncountably many parameter values that
correspond to no K3 family at all.

\section{Monodromy and the transcendental lattice}

The monodromy of $L$ acts on its solution space, i.e.\ on $T\otimes\C$, and
preserves the lattice $T$ and its form. Conversely, given only the operator,
one can hope to \emph{recover} $T$: compute the monodromy matrices around the
singular points numerically, find the invariant symmetric bilinear form, and
find the lattice generated by the orbit of a suitable vector. This is a
program rather than a theorem, and its output is only as good as its inputs:
the monodromy is computed numerically \Nstatus, the recognition of its
entries as rationals is a recognition step, and the identification of the
resulting rank-3 lattice with $T(X_z)$ requires knowing that the family is the
one the operator came from. Chapter~\ref{ch:frontier} reports one such
computation with every step labelled.

\section{Exercises}

\begin{exercise}
Compute $\Sym^2$ of $\theta^2-z(\theta+\tfrac12)^2$ (the hypergeometric
operator of $K(k)$) and identify the result as the Picard--Fuchs operator of
the Legendre family's symmetric square.
\end{exercise}

\begin{exercise}[$\star$]
Show that if $y_1,y_2$ solve $\theta^2y+p_1\theta y+p_0y=0$, then
$u=y_1y_2$ satisfies a third-order equation, and that $y_1^2,y_1y_2,y_2^2$
satisfy the quadratic relation $(y_1y_2)^2=y_1^2\cdot y_2^2$. This is the
``nondegenerate quadric'' of Doran's proof.
\end{exercise}

\begin{exercise}
Verify from Theorem~\ref{thm:cooper} that $s_7(1)=4$, $s_7(2)=48$,
$s_7(3)=760$ (with $s_7(0)=1$), and that $4\mid s_7(n)$ for $1\le n\le 10$.
(The general statement $4\mid s_7(n)$ for all $n\ge1$ is kernel-checked in the
program, \artifact{four\_dvd\_s7}.)
\end{exercise}

\section*{Chapter note}

Doran's theorem was read from a hash-pinned copy of~\cite{Doran2000} (its
Thm.~5.13, Cor.~5.8, and \S6). Zagier's and Cooper's classifications are
quoted from the cited papers; the parameter quadruples are those encoded and
golden-tested in the program's Lean development against the first terms of
the sequences as printed in the sources.
